% from ega4-2-fr.tex
\oldpage[IV-2]{63}

points maximaux de $X\otimes_k K$ sont ceux de la fibre
$p^{-1}(x)=\operatorname{Spec}(k(x)\otimes_k K)$ ((2.3.4) et
$(0_{\mathrm I},2.1.8)$)~; dire que cette fibre est irréductible pour toute
extension finie séparable $K$ de $k$ revient à dire que $k(y)$ est une
extension primaire de $k$, en vertu de (4.3.2). Inversement, si $c)$ est
vérifiée, le même raisonnement (compte tenu de (4.3.2)) montre que pour toute
extension $K$ de $k$, $X\otimes_k K$ est irréductible, donc $c)$ implique
$a)$.
\end{proof}

\begin{corollary}[4.5.10]
\label{IV.2.4.5.10-fr}
Soient $X$ un $k$-préschéma irréductible, $x$ son point générique, $k'$ la
fermeture algébrique séparable de $k$ dans $k(x)$, $k''$ une extension
galoisienne de $k$ (de degré fini ou non) contenant $k'$. On suppose $k'$ fini
sur $k$~; alors les composantes irréductibles de $X\otimes_k k''$ sont
géométriquement irréductibles~; leur nombre est égal à $[k':k]$ et est aussi le
nombre géométrique de composantes irréductibles de $X$.
\end{corollary}

\begin{proof}
Comme dans (4.5.9), on est ramené à considérer les points maximaux de la fibre
$\operatorname{Spec}(k(x)\otimes_k k'')$~; en vertu de (4.3.4), ils sont au
nombre de $[k':k]$ et les corps résiduels en ces points (qui sont aussi ceux de
$X\otimes_k k''$) sont des extensions primaires de $k''$, ce qui, compte tenu
de (4.5.9), prouve le corollaire.
\end{proof}

\begin{corollary}[4.5.11]
\EGAReaderTarget{ega4.c04.cor.4-5-11}{IV.2.4.5.11-fr}
Soit $X$ un $k$-préschéma~; supposons que $X$ n'ait qu'un nombre fini de
points maximaux $x_i$ $(1\leqslant i\leqslant r)$ et que pour chaque $i$, la
fermeture algébrique séparable $k'_i$ de $k$ dans $k(x_i)$ soit de degré fini
sur $k$. Alors il existe une extension finie séparable $L$ de $k$ telle que les
composantes irréductibles de $X\otimes_k L$ soient géométriquement
irréductibles~; leur nombre, égal à $\sum_i[k'_i:k]$, est le nombre géométrique
des composantes irréductibles de $X$.
\end{corollary}

\begin{proof}
Les points maximaux de $X\otimes_k L$ sont ceux des diverses fibres
$\operatorname{Spec}(k(x_i)\otimes_k L)$ en vertu de (2.3.4) et
$(0_{\mathrm I},2.1.8)$ appliqué aux composantes irréductibles de $X$ et à
leurs images réciproques dans $X\otimes_k L$. Il suffit donc de prendre pour
$L$ une extension galoisienne finie de $k$ contenant tous les $k'_i$ et
d'appliquer (4.5.10).
\end{proof}

On notera que (4.5.11) s'applique en particulier à tout $k$-préschéma $X$ de
type fini sur $k$, car $X$ est alors noethérien, donc n'a qu'un nombre fini de
composantes irréductibles, et chacun des $k(x_i)$ est une extension de type
fini de $k$, donc il en est de même de la fermeture algébrique de $k$ dans
$k(x_i)$ ([1], p.~6-06, lemme~5).

\begin{remark}[4.5.12]
\label{IV.2.4.5.12-fr}
\begin{enumerate}
\item[(i)] Les notions définies dans (4.5.2) dépendent du corps de base $k$.
Pour abréger, et lorsqu'il sera nécessaire de mentionner le corps de base, on
pourra dire «~$k$-irréductible~» (resp. «~$k$-connexe~») au lieu de
«~géométriquement irréductible (resp. connexe) relativement à $k$~». On notera
que cette terminologie, tout à fait naturelle dans le contexte des schémas,
est exactement à l'opposé de celle de Weil~: chez cet auteur,
«~$k$-irréductible~» (resp. «~$k$-connexe~», «~$k$-normal~», etc.) désigne une
propriété \emph{intrinsèque} d'un préschéma $X$, indépendante du corps $k$ pris
comme «~corps de base~» (autrement dit, du morphisme
$X\to\operatorname{Spec}(k)$ choisi)~; nous exprimons ces propriétés en disant
simplement que $X$ est irréductible (resp. connexe, normal, etc.). Weil dit
d'autre part «~absolument irréductible~», «~absolument connexe~», «~absolument
normal~», etc., pour les notions correspondantes relatives au préschéma
$X\otimes_k\Omega$, où $\Omega$ est une extension algébriquement close
convenable de $k$. Cette terminologie s'explique par le point de vue différent
auquel se place cet auteur~: pour lui, une «~variété algébrique~» $V$ est
donnée d'abord comme un objet

